\documentclass[11pt,a4paper]{article}
\usepackage[utf8]{inputenc}\usepackage[T1]{fontenc}
\usepackage[margin=1in]{geometry}
\usepackage{amsmath,amssymb,booktabs,graphicx,enumitem,xcolor,parskip}
\usepackage[hidelinks]{hyperref}
\usepackage{titlesec}
\definecolor{qnavy}{RGB}{20,40,75}\definecolor{qgrey}{RGB}{90,90,90}\definecolor{qred}{RGB}{150,30,30}
\titleformat{\section}{\large\bfseries\color{qnavy}}{\thesection}{0.6em}{}
\titleformat{\subsection}{\bfseries\color{qnavy}}{\thesubsection}{0.5em}{}

\title{\bfseries\color{qnavy} Pre-registration by an isolated agent: what it protects against,
and what it does not}

\author{Reza Nehzati, Ph.D.\\[2pt]
{\small\itshape VMC MAR COM Inc.\ DBA HeyDonto, Knoxville, TN, United States}}
\date{16 August 2026}

\begin{document}
\maketitle
\thispagestyle{empty}

\begin{abstract}
\noindent
Language-model agents can now write and execute clinical data analyses. The obvious hazard is not
incompetence but flexibility: a system that sees the data before fixing its plan can produce a
defensible-looking finding from almost any dataset. We test a protocol designed to remove that
flexibility, on three sealed oncology cohorts. A generator agent receives only schema and
documentation --- form names, column names, row counts, a data dictionary, and no cell value from
any outcome field --- and must emit a complete pre-registered plan: hypotheses, positive controls
with pre-declared failure consequences, integrity gates, an exhaustive exclusion ledger, a power
statement, and mechanical verdict rules. The plan is SHA-256 hashed before any outcome is read,
and the exposure table is frozen and hashed before outcomes are merged, so post-outcome tuning is
mechanically detectable rather than merely discouraged.
Across three sealed cohorts the system produced three distinct behaviours. On the first
(247 stage-IV NSCLC patients treated with anti-PD-(L)1) it returned POWER-LIMITED-NULL rather
than a finding: the pre-specified multimodal comparison was
structurally underpowered at $n = 73$, and the plan's own rule required saying so. On the second
(ACRIN 6668, 166 analysed patients, 125 events), a structurally isolated generator produced a plan
that recovered the trial's designed endpoint --- post-treatment FDG peak SUV predicting shorter
survival, HR 1.202 per doubling, permutation $p = 0.0005$ --- having identified four
analysis-destroying traps from schema alone, including immortal-time bias and a vital-status code
that silently undercounts deaths.
On the third (46 anti-PD-1 patients with imaging but no released outcome) it declined outright,
returning \texttt{feasible: false} with the reasoning that outcome missingness was 100\% and the
exposure was a \emph{constant} --- every patient received the drug that defines the collection, so
no contrast exists --- and it identified a covariate leak in the metadata we had supplied without
noticing it ourselves.
The protocol's limits are as instructive as its successes. In the first cohort a pre-registered
hard gate failed and was reinterpreted after the fact. In the second, an implementation error let
a control's hazard ratio be reported as the primary result while all eight gates passed. We
conclude that pre-registration constrains \emph{analytic} discretion effectively and
\emph{implementation} correctness not at all, and that the two require different defences. Two further
gaps are now visible: neither defence prevents a small-sample estimate being reported as though it
were stable, and neither prevents a plan from measuring a confounded quantity with complete
procedural correctness. The last is the most serious, because the protocol's own machinery reports
success throughout.
\end{abstract}

\section{Introduction}

An agent that can read a data dictionary, write analysis code and execute it can also, without any
intent to deceive, arrive at a result by trying things until something works. The number of
defensible-looking choices in a survival analysis --- time origin, exclusion rules, transformation,
cutpoint, covariate set, handling of ambiguous status codes --- is large enough that a
sufficiently persistent search will find a significant one.

Pre-registration is the standard answer, and it transfers to this setting with one modification
that matters: the plan can be written by a system that has provably not seen the outcomes, and the
ordering can be enforced by a hash rather than by a declaration.

We ask two questions. Does such a system produce a plan a competent analyst would accept? And
does it refuse to manufacture findings the data cannot support? We report three cohorts --- one
confirmation, one power-limited null, one outright refusal --- and we report the protocol's
failures at equal length.

\section{Protocol}

\begin{enumerate}[leftmargin=1.4em,itemsep=3pt]
\item \textbf{Recon package.} The generator receives a single JSON file: every form's name, column
      names and row counts; the trial's data dictionary; public facts about the cohort; and
      join counts. It contains no cell value from any outcome, measurement or clinical field.
\item \textbf{Isolation.} In the second and third cohorts the generator was a separate agent whose
      readable surface was a directory containing only that file. The location of the actual data was never
      disclosed to it.
\item \textbf{Plan.} The generator emits hypotheses with named fields, positive controls each with
      an explicit pre-declared failure consequence, numbered integrity gates with expected values,
      an exhaustive exclusion ledger, a power statement computed before execution, multiplicity
      handling, mechanical verdict rules, and a list of conclusions it will refuse to draw.
\item \textbf{Hash before labels.} The coordinator SHA-256 hashes the plan and records the
      timestamp before reading any outcome value.
\item \textbf{Staged execution with an exposure freeze.} Structural gates run first on keys and
      counts only. The exposure table is then built, written and hashed. Only then may outcome
      forms be opened --- enforced in the second cohort by a loader that raises if asked for an
      outcome form before the freeze. A later gate re-hashes the exposure table and compares,
      making post-outcome modification of the predictor mechanically detectable.
\item \textbf{Mechanical verdict.} Verdict rules are evaluated from the numbers, in a declared
      precedence order, with no discretion at the final step.
\end{enumerate}

\section{Cohort 1 --- refusal under low power}

247 stage-IV NSCLC patients treated with anti-PD-(L)1, with RECIST response, survival, PD-L1
immunohistochemistry slides, targeted panel sequencing and CT radiomics. The original authors'
analysis code was placed under a self-imposed firewall for the whole arm.

\subsection{What the system did}

Eleven gates were declared and ten passed, including hash-before-label ordering, exclusion
accounting that reconciles per endpoint, absence of NaN or constant prediction vectors with 100\%
out-of-fold coverage, and permutation nulls centred at 0.494--0.495 across every model family.
That null centring is what makes the accompanying $p$-values worth reading.

The pre-declared positive control --- PD-L1 tumour proportion score, \emph{on its own} ---
predicted response at AUROC \textbf{0.747} (permutation $p = 1.0\times10^{-4}$, $n = 246$),
clearing its declared bar of AUROC $>0.55$ at $p<0.05$. The multivariable clinical baseline that
the primary comparison is measured against --- PD-L1 plus tumour mutational burden plus clinical
variables --- reached AUROC \textbf{0.745} on the same 246 patients (0.742 as the mean over
cross-validation repeats).

\textcolor{qred}{\textbf{Corrected 2026-08-17.}} This paragraph originally attributed the 0.747 to
the multivariable baseline. It is the single-variable positive control; the baseline model is
0.745. The number was taken from the evidence but assigned to the wrong model, and the error
survived a verification pass that compared values without checking which field produced them. It
does not affect the primary comparison below, which is a paired contrast against the multivariable
baseline on the complete-case subset. Found by the scoped re-verification described in
\texttt{papers/verify\_claims.py}.

The pre-specified primary comparison asked whether adding imaging and pathology to that baseline
improved it. On the 73 patients with every modality present, the change was $-0.036$
(95\% CI $-0.220$ to $+0.150$, one-sided $p = 0.638$). The plan's rule mapped this to
\textbf{POWER-LIMITED-NULL}: a bootstrap interval half-width of 0.185 cannot distinguish a useful
gain from a harmful one, so no finding was claimed in either direction.

Three further modality additions, each pre-declared, were all negative and none individually
significant: radiomics $-0.016$, pathology-global $-0.019$, pathology-focal $-0.056$. Four
negative estimates out of four is more informative than any single interval, and the honest
reading is that at these sample sizes added parameters cost more than they return.

Holm correction across the model family left two survivors of six. The system reported the four
that did not survive as not claimable, including one at raw $p = 0.018$.

\subsection{An unexpected observation}

Embeddings from a general-purpose pathology foundation model predicted a PD-L1 tumour proportion
score of $\geq 50\%$ at AUROC \textbf{0.870} (95\% CI 0.811--0.922, $n = 163$, 67 positive) with no
immunohistochemistry-specific training. This was a competence control, not a hypothesis; we report
it because it was surprising, and note that a second definition of the same quantity gives 0.856,
so the figure is definition-dependent.

\subsection{The protocol's first failure}

A pre-registered \emph{hard} gate failed. The control asked whether segmentation-derived tumour
volume predicted survival, with the stated basis ``proves I can read the imaging''. Observed:
C-index \textbf{0.5035}, Wald $p = 0.698$, on 187 patients with 111 events. Flat.

The asymmetry in the plan is the problem. Its pathology control carried an explicit remedy written
in advance --- ``failure $\Rightarrow$ report pipeline not competent, not evidence of absence''.
The imaging control carried none: it was a bare pass/fail gate. So when it failed there was no
pre-registered instruction, and a radiomics model built on that same imaging was one of only two
Holm survivors.

A defence is available and reasonable --- the control tested whether target-lesion segmentations
approximate \emph{total} tumour burden, a different and questionable premise, and three other
controls passed. But it is a \textbf{post-hoc reinterpretation of a pre-registered control}, and
calling it anything else would defeat the purpose of the exercise. The instructive part is that
pre-registration surfaced a control whose basis was mis-specified, which is what pre-registration
is for.

\section{Cohort 2 --- confirmation under isolation}

ACRIN 6668 / RTOG 0235: FDG-PET before and after definitive chemoradiotherapy in stage III NSCLC,
released with clinical data under CC BY 3.0. Here the generator was structurally isolated as
described in \S2, and produced a 493-line plan hashed at
\texttt{752b7462fce56899\ldots cb74e6770} before any outcome was read.

\subsection{What it caught from schema alone}

\begin{itemize}[leftmargin=1.3em,itemsep=2pt]
\item \textbf{Immortal time.} The exposure is measured 8--20 weeks after treatment, so measuring
      survival from registration would guarantee survival to exposure. It set the time origin at
      the post-treatment scan and made $\min(T)\geq 0$ a gate.
\item \textbf{A vital-status code meaning ``dead, date of death unknown''.} The natural
      \texttt{event = (status == 2)} silently undercounts deaths. It treated the code as an event
      and made the opposite coding a \emph{condition} of confirming.
\item \textbf{Cutoff forking.} One form carries 1{,}185 rows keyed by an SUV inclusion threshold ---
      roughly five analyses per patient. It locked to a single threshold and forbade comparison
      across them.
\item \textbf{An undocumented uniqueness assumption.} Its primary exposure required one row per
      patient in a form where the recon package never said so. It demanded verification and
      pre-specified a duplicate-resolution rule. The rule fired: 244 rows over 243 patients.
\end{itemize}

It also pre-declared fifteen conclusions it would refuse to draw, and forbade cutpoint searches,
subgroup fishing, and endpoint substitution.

\subsection{Result}

\begin{center}
\begin{tabular}{@{}llr@{}}
\toprule
Primary: OS on log$_2$ post-treatment peak SUV & HR per doubling & \textbf{1.202} \\
& 95\% CI & 1.100--1.315 \\
& permutation $p$ ($B = 2000$) & \textbf{0.0005} \\
& analysed / events & 166 / 125 \\
\midrule
PC0 exposure validity (outcome-free) & post below pre & PASS (94.4\%) \\
PC1 progression $\rightarrow$ death & HR 2.463 & PASS ($p = 3.8\times10^{-5}$) \\
PC2 performance status & HR 1.382 & PASS-WEAK ($p = 0.072$) \\
PC3 weight loss $\geq 5\%$ & 26 exposed & UNINFORMATIVE-BY-N \\
\bottomrule
\end{tabular}
\end{center}

Ten gates passed, the exclusion ledger reconciled exactly ($251 = 166 + 85$), the permutation null
was centred and calibrated (mean $z$ 0.034, SD 0.982, rejection rate 0.042), and the exposure hash
was unchanged from the Stage-1 freeze. All seven pre-declared robustness tests agreed in direction
and survived Holm correction, including the adjusted model (1.214), a fixed threshold with no
search permitted (1.789), and the opposite death-coding (1.202).

\textbf{Verdict: CONFIRMED} --- the trial's designed endpoint, recovered by a plan written in
ignorance of the data's contents.

Two controls did not fully corroborate. Performance status was directionally right but not
significant; weight loss was declared uninformative because 26 patients met the exposure threshold
against a pre-declared minimum of 30. Both had been \emph{anticipated}: trial eligibility screens
performance status, so range restriction was given its own outcome branch in advance rather than
explained afterwards.

The plan also forced an uncomfortable number into the open. \textbf{106 of 180 exposure values
were at or below the floor} --- complete metabolic response --- so 59\% of the cohort sits at one
tied value and the association is carried by the minority with residual uptake. A plan written
after seeing the data would have had every opportunity not to mention this.

\subsection{The protocol's second failure, and the more important one}

An earlier execution reported the primary as HR 1.382 with 95\% CI 1.050--1.124: a point estimate
outside its own confidence interval, which is impossible. Inside the control blocks the coordinator
had reused two variable names, clobbering the primary's values before the verdict logic read them.
The verdict was therefore briefly evaluated using a \emph{control's} hazard ratio.

\textbf{All eight of the plan's gates passed while this was true.} The gates checked recombination,
joins, cutoff inflation, leakage, exclusion accounting, permutation calibration, exposure coverage
and outcome completeness. None of them could see that the number being reported was not the number
that had been computed. What caught it was an internally impossible pair of values, noticed by a
human reading the output.

The conclusion survives correction. That it was momentarily correct by accident is the point.

\section{Cohort 3 --- refusal when no analysis exists}

The two cohorts above both had endpoints. A protocol that only ever produces analyses has not
demonstrated the second half of the claim, so we gave the isolated generator a cohort with none:
46 non-small-cell lung cancer patients treated with anti-PD-1, 677 imaging series across CT, PET
and three ``other'' series, pre-treatment and one follow-up timepoint. The archive released
images and a retrieval manifest and nothing else --- no response assessment, no survival, no
demographics, no stage, no lesion annotation.

It was asked whether a valid pre-registered analysis was possible and to write one if so. It
returned \texttt{feasible: false} in 567 lines.

\subsection{Its reasoning}

Two defects, each independently fatal. Outcome missingness is 46 of 46: not a missing-data
problem, because there is no observed stratum to model from and no estimand can be defined. And
\textbf{the exposure is a constant} --- the collection is \emph{defined} by every patient having
received anti-PD-1, so there is no contrast, no comparator, and no released regimen or line
variation to recover one. The second point is the sharper of the two and we had not articulated
it ourselves.

It nonetheless computed the power ceiling in advance, so that a later ``repaired'' version could
not quietly overreach: at $n = 46$ with roughly 32 events the minimum detectable hazard ratio is
about 2.7, which puts signature \emph{discovery} permanently out of reach and admits only a
confirmatory test of a pre-frozen external signature.

\subsection{The shortcuts it named and refused}

Most valuable is what it identified as tempting rather than merely absent. The gravest is a
pseudo-RECIST endpoint built from the single follow-up scan: a volumetric change looks like a
reader-free quantitative outcome and would plausibly survive review. It is disqualified
\emph{specifically in this cohort} because under anti-PD-1 early enlargement can reflect immune
infiltration rather than growth, so the surrogate misclassifies in a direction correlated with
the biology under study --- bias that points at the hypothesis rather than adding noise.

It also flagged the three ``other'' series as bait for optical character recognition, since dose
reports and console captures are where rendered text survives, and noted this would both cover
too few patients and attempt to defeat de-identification; free-text protocol headers as a source
of grep-able outcome labels such as ``restaging'' or ``progression'', which are scheduling
artefacts with no denominator; and a maximum of 36 series per patient as reconstruction
redundancy rather than 36 timepoints.

\textbf{And it found a leak we had introduced.} The recon package listed the most common protocol
names, among them \texttt{1.2 BMI 0-25 FDG Body} and \texttt{BMI\_25to40\_PETCTBody\_CBM}. These
leak a coarse body-habitus covariate through weight-stratified PET protocols, with an
obesity-paradox literature hook ready to receive it. We had included those strings without
noticing. It declined to use them --- protocol-selection labels are not measurements, the bands do
not partition, they are observed only in PET-imaged patients, and they reconstruct a variable the
custodian chose not to release.

\subsection{What it offered instead}

Rather than pad the refusal, it proposed an outcome-free study explicitly fenced as a different
question --- about the instrument rather than the patient: a header-only harmonisability audit
with a pre-declared rule that if the largest homogeneous acquisition stratum falls below 15
patients, ``this collection cannot support radiomics'' \emph{is} the published result. Its
extension turns the reconstruction redundancy into a repeatability estimate using automatic
non-lesion reference organs, which removes the reader-discretion problem outright rather than
mitigating it.

It closed by conditioning its own verdict: the refusal rests on the recon package, and does not
survive a custodian demonstrating that a patient-level clinical supplement exists.

\section{Four occasions when the discipline actually cost something}

The three cohorts above show a protocol behaving well. They do not show it \emph{binding}, because in
none of them did a pre-declared rule prevent a conclusion we wanted. Four later runs did, and they
are worth stating because a constraint that never bites is decoration. Two concern the protocol's
mechanics; two concern its limits, and those are the more useful.

\subsection{A threshold that would have been convenient to move}

A five-encoder survey of site-leakage in whole-slide features reused verdict thresholds fixed and
hashed in an earlier run, before that run's own results were read. One encoder returned a mean
methylation inflation of $+0.0173$ where the other three histology encoders gave $+0.0606$ to
$+0.0849$.

Had the deciding rule been written that day it would naturally have been a magnitude bar --- and at
any bar near 0.05 that encoder fails, turning a clean four-for-four result into a mixed one. The rule
that actually existed keyed on \emph{how many of six targets} inflated rather than by how much, so it
fired, and the fivefold spread went into the limitations where it belongs. We could not have chosen
either framing honestly after seeing the number. The pre-registration made the choice for us, and it
happened to choose the stricter-in-kind option.

\subsection{A pattern we are not allowed to claim}

The same survey produced its most interesting observation by accident: relative inflation falls as
site-disjoint capability rises, Spearman $\rho = -0.80$ across five encoders. It looks like a law when
plotted.

It is not reported as a finding, because it was noticed \emph{after} all five arms were read, and at
$n = 5$ no rank correlation can reach $p<0.05$ by that test --- the observed $p$ is 0.104. Testing a
hypothesis on the data that suggested it is precisely the flexibility this paper is about. It is
recorded in the run's own output under a field whose value is the string
\texttt{POST-HOC OBSERVATION, NOT A TESTED FINDING}, and what it licenses is a pre-registration on
encoders not yet used, not a conclusion.

This is the uncomfortable case. The most publishable observation in that run is the one we are least
entitled to, and the protocol offers no way to have it.

\subsection{And a gate that caught a statistic rather than a bug}

A later run re-measured an external cohort separation whose value had been reported as 0.8768. A
reproduction gate pinned that number as a hard expectation. Re-run with a different random seed the
same quantity came out 0.9240, and the gate halted the run.

Nothing was broken. The seed sets the cross-validation split, and on 76 samples with roughly fifteen
per fold a few reassignments move balanced accuracy several points; measured over 25 splits the value
is $0.9176 \pm 0.0247$. The originally reported 0.8768 was the pre-declared seed --- so not chosen to
flatter anything --- but it sat 1.6 standard deviations below its own mean and had been quoted as a
fixed quantity. The earlier report's conclusion survived; its headline contrast did not.

The gate was written to detect a coding difference and detected a statistical one. This is a third
failure mode, distinct from both of the above: not analytic flexibility, not implementation error, but
\emph{treating a noisy estimate as a point}. It suggests a fourth defence --- report distributions
over resampling wherever $n$ is small --- which none of our pre-registrations required.

\subsection{A plan that measured the wrong thing, competently}

The protocol was applied to a question about drug resistance: in 187 lung cancer cell lines with
450k methylation and matched dose--response, are lines resistant to standard therapy differentially
sensitive to chromatin-targeting drugs? The plan was hashed before any half-maximal inhibitory
concentration was joined to any methylation value.

It returned a clean, significant, pre-registered answer: resistant lines are \emph{cross}-resistant,
Spearman $\rho = +0.651$, permutation $p = 0.0005$, $n = 187$. Under the declared rule that is
THESIS\_CONTRADICTED.

Then a control we added afterwards destroyed it. Both indices were mean $z$-scored ln\,IC$_{50}$, and
cell lines differ enormously in \emph{general} drug sensitivity --- some are resistant to nearly
everything. Partialling out an index built from 227 drugs in neither set takes $+0.651$ to
$\mathbf{-0.160}$. The sign reverses.

We are not entitled to swap one number for the other: the post-hoc figure is not a pre-registered
test and cannot replace one. What the pair licenses is the conclusion that \emph{the design was
confounded and answers nothing}.

\textbf{This is the sharpest limitation we have found.} A correctly executed pre-registration, whose
every gate passed, produced a publishable-looking and wrong conclusion. Had we not been suspicious of
how clean $+0.651$ looked, the report would have stated a manuscript's central premise contradicted
at $p = 0.0005$ in 187 cell lines. \textbf{Pre-registration constrains analytic flexibility; it does
not confer construct validity.} A hashed plan that measures the wrong quantity measures the wrong
quantity with excellent discipline, and nothing in the protocol detects that.

\subsection{A two-sided rule that cost us the headline}

The same protocol was then applied to a real intervention rather than a correlation. SMARCA4 is the
ATPase of the SWI/SNF chromatin-remodelling complex, so knocking it down in an osimertinib-resistant
lung cancer line is epigenetic modulation of a resistant state. The plan asked whether that reverses
the resistance transcriptional programme, and was hashed after reading only the column headers of the
expression matrix.

The knockdown worked --- SMARCA4 $\log_2$FC $-2.37$ against a pre-declared bar of $-0.5$, checked as a
gate. Genes up in resistance then fell below the permutation null exactly as the thesis predicts
($p = 0.0005$). Under a one-sided rule that is a positive result.

The rule declared in advance was two-sided: reversal required resistance-up genes to fall \emph{and}
resistance-down genes to rise. The down set fell as well, so the mechanical verdict was PARTIAL. The
two-sided form is the correct one --- a treatment that lowers everything associated with resistance is
suppressing a correlated set, not reprogramming a state, and only the up-down pattern separates
those. Writing it cost nothing in advance and cost the headline when it fired.

A replication arm carried in the same plan, with no threshold declared, then inverted every sign in a
second model. One model showing partial reversal and another showing the opposite is not a finding.
Without that arm the report would have claimed partial reversal at $p = 0.0005$.

\section{Discussion}

\subsection{What pre-registration bought}

Concretely, in these three cohorts: a null was reported as a null with its power stated rather than
as an absence of effect; four of six models were declared not claimable after multiplicity
correction, including one at raw $p = 0.018$; a control was declared uninformative rather than
counted as a pass; an inconvenient distributional fact about 59\% of the exposure was reported;
four analysis-destroying traps were identified before any data was touched; and a fixed cutpoint
was used where a search would have been available and rewarding; and an entire cohort was
declined rather than analysed with a surrogate endpoint that would likely have passed review.

None of that requires trusting the analyst. The hash ordering and the exposure freeze make the
central claims checkable by a third party from the artefacts alone.

\subsection{What it did not buy}

\textbf{It does not make a plan correct.} The first cohort's imaging control was mis-specified: its
stated basis did not match what it actually tested. Pre-registration fixed that mistake in place
where it could be seen, which is better than the alternative, but it did not prevent it.

\textbf{It does not protect against implementation error.} This is the finding we would most want
carried forward. A pre-registration constrains which analysis you are allowed to run. It says
nothing about whether the code you wrote runs that analysis. In the second cohort every gate passed
while the reported primary was a different quantity from the computed one, and the discrepancy was
visible only because two numbers were mutually inconsistent.

The remedy is not more pre-registration. It is a distinct class of check: internal-coherence
assertions on the outputs themselves --- a point estimate inside its interval, a tally summing to
its total, a hash matching its earlier value. We added the first of these after the fact. A mature
protocol should declare them alongside the analytic gates, because \textbf{analytic discretion and
implementation correctness are different failure modes and neither defence covers the other}.

\subsection{Limitations}

\textbf{Three cohorts and four later runs.} One confirmation, one power-limited null, one outright
refusal, plus four occasions where a pre-declared rule bound against our interest --- including one where a
correctly executed plan produced a confidently wrong answer. That is a
demonstration, not an evaluation, and it cannot support a claim about how often such a system
behaves well.

\textbf{Only the second generator was structurally isolated.} In the first, one agent restrained
itself, and a self-imposed firewall is a weaker guarantee than an inaccessible one.

\textbf{The second cohort's analysis was tabular.} It used core-laboratory PET assessments, not the
underlying image volumes, so the claim concerns an unseen cohort and data shape rather than an
unseen imaging pipeline.

\textbf{Neither result is an independent replication.} The second cohort's finding is the trial's
own published conclusion, obtained from the same data. The claim is about process, not biology.

\textbf{One coverage gate passed at exactly its threshold} (180 against $\geq$ 180). One patient
fewer would have triggered a declared fallback and changed the exposure source, which illustrates
how much a pre-registered threshold can hinge on an arbitrary boundary.

\textbf{The generator's isolation rests partly on its own report} of which files it read, alongside
the structural constraint that the data's location was never disclosed and its readable directory
contained one file.

\section*{Data and code availability}

Both cohorts are public: the first via Synapse, the second via The Cancer Imaging Archive under
CC BY 3.0. The hashed plan, its hash and timestamp, the recon package given to the generator, the
executor, the frozen exposure table with its hash, gate records, exclusion ledgers, all permutation
statistics, and out-of-fold predictions are retained and available on request.

\section*{Competing interests}

The author is affiliated with a commercial entity developing computational pathology methods.

\end{document}
